\section{\texorpdfstring{Násobení a dělení}{Nasobeni a deleni}}

% =====================================================================================================================
  \begin{frame}
    \frametitle{Násobení ve 2. doplňku}
    \boldmath
    
    \begin{itemize}
      \item Násobení v ostatních soustavách funguje stejně jako v desítkové soustavě - násobení pod sebou,
      \item je extrémně jednoduché pro binární soustavu, protože opisujeme jen nenulové řádky,
      \item pokud rozšiřujeme správně, pak kódování ve 2. doplňku funguje při násobení zcela automaticky,
      \item výsledek násobení může mít tolik řádů, kolik je součet počtu řádů obou členů.
    \end{itemize}
    
      \textbf{Příklad:}    
      \begin{align*}
        (3)_{10} \cdot (3)_{10} =& (9)_{10} \\
        (11)_{2} \cdot (11)_{2} =& (1001)_{2}
      \end{align*}
      
      Oba členy mají 2 místa, výsledek má 4.
	\end{frame}
  
% =====================================================================================================================
  \begin{frame}
    \frametitle{Násobení ve 2. doplňku - pod sebou, kladné členy}
    \boldmath
    
    \textbf{Příklad:} $6\cdot 9 = 54$
    
      \vspace{3mm}
      1. člen: $(110)_{2}$, po zavedení znaménka $(0110)_{2}$
      
      2. člen: $(1001)_{2}$ , po zavedení znaménka $(01001)_{2}$
      
      $\Rightarrow$ oba členy rozšíříme nejméně na 9 bitů - zde 10 bitů
      \vspace{10mm}
      \begin{center}
        \begin{tabular}{c | c c c c c c c c c c}
          $\longleftarrow$ \textcolor{red}{0} & \textcolor{blue}{0} & \textcolor{blue}{0} & \textcolor{blue}{0} & \textcolor{blue}{0} & \textcolor{blue}{0} & 0 & 1 & 0 & 0 & 1 \\
          $\longleftarrow$ \textcolor{red}{0} & \textcolor{blue}{0} & \textcolor{blue}{0} & \textcolor{blue}{0} & \textcolor{blue}{0} & \textcolor{blue}{0} & 0 & 0 & 1 & 1 & 0 \\
          \hline
          $\longleftarrow$ \textcolor{red}{0} & \textcolor{blue}{0} & \textcolor{blue}{0} & \textcolor{blue}{0} & \textcolor{blue}{0} & \textcolor{blue}{0} & 0 & 0 & 0 & 0 & 0 \\
          $\longleftarrow$ \textcolor{red}{0} & \textcolor{blue}{0} & \textcolor{blue}{0} & \textcolor{blue}{0} & \textcolor{blue}{0} & 0 & 1 & 0 & 0 & 1 &  \\
          $\longleftarrow$ \textcolor{red}{0} & \textcolor{blue}{0} & \textcolor{blue}{0} & \textcolor{blue}{0} & 0 & 1 & 0 & 0 & 1 &  &  \\
          \hline
          $\longleftarrow$ \textcolor{red}{0} & 0 & 0 & 0 & 0 & 1 & 1 & 0 & 1 & 1 & 0
        \end{tabular}
      \end{center}
    
	\end{frame}

% =====================================================================================================================
  \begin{frame}
    \frametitle{Násobení ve 2. doplňku - pod sebou, záporný člen}
    \boldmath
    
    \textbf{Příklad:} $6\cdot (-9) = -54$
    
      \vspace{3mm}
      1. člen: $(110)_{2}$, po zavedení znaménka $(0110)_{2}$
      
      2. člen: $(10111)_{2}$
      
      $\Rightarrow$ oba členy rozšíříme nejméně na 9 bitů - zde 10 bitů
      \vspace{10mm}
      \begin{center}
        \begin{tabular}{c | c c c c c c c c c c}
          $\longleftarrow$ \textcolor{red}{1} & \textcolor{blue}{1} & \textcolor{blue}{1} & \textcolor{blue}{1} & \textcolor{blue}{1} & \textcolor{blue}{1} & 1 & 0 & 1 & 1 & 1 \\
          $\longleftarrow$ \textcolor{red}{0} & \textcolor{blue}{0} & \textcolor{blue}{0} & \textcolor{blue}{0} & \textcolor{blue}{0} & \textcolor{blue}{0} & 0 & 0 & 1 & 1 & 0 \\
          \hline
          $\longleftarrow$ \textcolor{red}{0} & \textcolor{blue}{0} & \textcolor{blue}{0} & \textcolor{blue}{0} & \textcolor{blue}{0} & \textcolor{blue}{0} & 0 & 0 & 0 & 0 & 0 \\
          $\longleftarrow$ \textcolor{red}{1} & \textcolor{blue}{1} & \textcolor{blue}{1} & \textcolor{blue}{1} & \textcolor{blue}{1} & 1 & 0 & 1 & 1 & 1 &  \\
          $\longleftarrow$ \textcolor{red}{1} & \textcolor{blue}{1} & \textcolor{blue}{1} & \textcolor{blue}{1} & 1 & 0 & 1 & 1 & 1 &  &  \\
          \hline
          $\longleftarrow$ \textcolor{red}{1} & 1 & 1 & 1 & 1 & 0 & 0 & 1 & 0 & 1 & 0
        \end{tabular}
      \end{center}
    
	\end{frame}
  
% =====================================================================================================================
  \begin{frame}
    \frametitle{Dělení ve 2. doplňku}
    \boldmath
    
    \begin{itemize}
      \item Dělení v ostatních soustavách funguje podobně jako v desítkové soustavě,
      \item pro záporná čísla lze používat \textbf{metodu dělení s obnovou zbytku} nebo \textbf{převedení na absolutní hodnoty},
      \item existuje zde problém přetečení při dělení minima rozsahu hodnotou -1:
    \end{itemize}
    \vspace{5mm}
    
      \textbf{Příklad:} 
      
       8 bitový registr s hodnotou se znaménkem ve 2. doplňku.
        
        \begin{itemize}
          \item Rozsah hodnot -128 až 127,
          \item problematická operace je $-128/-1 = 128$, což je mimo rozsah.
        \end{itemize}
      
      $\Longrightarrow$ Musím rozšířit dělenec o 1 bit! Výsledek dělení je také širší o 1 bit.
	\end{frame}
 
% =====================================================================================================================
  \begin{frame}
    \frametitle{Dělení ve 2. doplňku - kladné členy}
    \boldmath
    \small
    \textbf{Příklad:} $15 / 3 = 5$
    
      1. člen: $(1111)_{2}$, po zavedení znaménka $(01111)_{2}$
      
      2. člen: $(11)_{2}$ , po zavedení znaménka $(011)_{2}$, záporná hodnota $(11101)_{2}$
      
      \vspace{3mm}
      \begin{center}
        \begin{tabular}{r || r | r r r r r r r r}
         \textcolor{blue}{\textbf{zbytek}} & & 001111 & : 11 =  \textcolor{red}{000}\textcolor{blue}{1}\textcolor{red}{0}\textcolor{blue}{1} \\
         posun dělitele a odečtení od zbytku & 1110 & 100000 & + \textcolor{blue}{001111}\\
          \hline
          \textcolor{red}{\textbf{záporný výsledek}} & 1110 & 101111 & \\
          posun dělitele a odečtení od zbytku & 1111 & 010000 & + \textcolor{blue}{001111}\\
          \hline
          \textcolor{red}{\textbf{záporný výsledek}} & 1111 & 011111 & \\
          posun dělitele a odečtení od zbytku & 1111 & 101000 & + \textcolor{blue}{001111}\\
          \hline
          \textcolor{red}{\textbf{záporný výsledek}} & 1111 & 110111 & \\
          posun dělitele a odečtení od zbytku & 1111 & 110100 & + \textcolor{blue}{001111}\\
          \hline
          \textcolor{blue}{\textbf{kladný výsledek, zbytek}} & 0000 & 000011 & \\
          posun dělitele a odečtení od zbytku & 1111 & 111010 & + \textcolor{blue}{000011}\\
          \hline
          \textcolor{red}{\textbf{záporný výsledek}} & 1111 & 111101 & \\
          posun dělitele a odečtení od zbytku & 1111 & 111101 & + \textcolor{blue}{000011}\\
          \hline
          \textcolor{blue}{\textbf{kladný výsledek, zbytek}} & 0000 & 000000 & \\
        \end{tabular}
      \end{center}
    
	\end{frame}